\documentclass[../NDCNNAI_main.tex]{subfiles}
\graphicspath{{\subfix{../fig/}}}

\begin{document}
\subsection{Предварительные сведения}
В классической теории нейронных сетей универсальное аппроксимационное свойство
(Universal Approximation Property, УАС) доказывается для конкретных архитектур
и, как правило, на компактных множествах. Однако современные модели обучения
выходят за рамки полносвязных сетей и включают операторные сети, сверточные и
структурированные архитектуры, а также аппроксимацию в бесконечномерных
функциональных пространствах. Это требует более абстрактного взгляда на УАС,
независимого от конкретной архитектуры. Наша цель --- ввести
мета-теоретический подход к универсальной аппроксимации, предложенный
А.~Крациосом, и зафиксировать используемые топологические и функциональные
понятия.

Поскольку универсальное аппроксимационное свойство формулируется как утверждение
о плотности, его корректная постановка требует явного выбора функционального
пространства и топологии на нём. В работе \cite{Kratsios2021} рассматриваются
различные классы пространств функций, возникающие в задачах обучения, включая
пространства непрерывных, интегрируемых и более общих функций.

Одним из базовых примеров является пространство непрерывных функций
\[
C(\Omega,\mathbb{R}^n) \coloneq \{ f \colon \Omega \to \mathbb{R}^n \mid f \text{ непрерывна} \}.
\]
где $\Omega \subset \mathbb{R}^d$ --- открытое или компактное множество.
Если $\Omega$ компактно, то естественной является норма супремума
$\|f\|_\infty = \sup_{x\in\Omega}\|f(x)\|$.
Однако при некомпактной области $\Omega$ равномерная сходимость на всём множестве
становится слишком жёстким требованием. В этом случае используется топология
\emph{равномерной сходимости на компактах}.

Для её задания фиксируется расширяющася система компактных подмножеств 
$(K_m)_{m\ge1}$ множества $\Omega$, то есть последовательность компактов,
удовлетворяющая $K_m \subset K_{m+1}$ и $\bigcup_{m=1}^\infty K_m = \Omega$.
На пространстве $C(\Omega,\mathbb{R}^n)$ вводится метрика
\[
d_{\mathrm{ucc}}(f,g)
=
\sum_{m=1}^{\infty}
2^{-m}
\frac{\|f-g\|_{\infty,K_m}}{1+\|f-g\|_{\infty,K_m}},
\qquad
\|h\|_{\infty,K} \coloneq \sup_{x\in K}\|h(x)\|.
\]
Данная метрика порождает топологию равномерной сходимости на компактах, которая
в данном контексте совпадает с компактно-открытой топологией
\cite[Example~1, Section~2.1]{Kratsios2021}.

Помимо пространств непрерывных функций, рассматриваются пространства
интегрируемых функций $L^p_\mu(\mathbb{R}^m,\mathbb{R}^n)$, где $\mu$ --- борелевская
мера, а расстояние индуцируется нормой
\[
\|f\|_{p,\mu} = \left(\int_{\mathbb{R}^m}\|f(x)\|^p\,d\mu(x)\right)^{1/p}.
\]
В отличие от $C(\Omega,\mathbb{R}^n)$, такие пространства являются банаховыми и
обладают линейной структурой, что существенно упрощает анализ операторов и
аппроксимационных свойств.

Ключевым наблюдением является то, что не все пространства, возникающие в задачах
машинного обучения, являются линейными. Примеры включают пространства форм,
графов, деревьев и вероятностных мер. Для унификации анализа Крациос использует
понятие \emph{точечного метрического пространства}
$(X,d_X,0_X)$, где фиксирован выделенный элемент $0_X$.

Для таких пространств вводится класс липшицевых отображений
$\mathrm{Lip}_0(X,Y)$, сохраняющих выделенную точку и не искажающих расстояния
сверх линейного множителя. Существенную роль играет конструкция
\emph{липшицева свободного пространства} $\mathcal{B}(X)$, позволяющая
линеаризовать нелинейное метрическое пространство. Существует каноническое
отображение
\[
\delta_X\colon X \to \mathcal{B}(X),
\]
которое является изометрией и обладает универсальным свойством:
любое липшицево отображение $f\colon X \to Y$ в банахово пространство $Y$ допускает
единственное линейное продолжение через $\delta_X$
\cite[Example~3, Section~2.1]{Kratsios2021}.

Наконец, большинство функциональных пространств, рассматриваемых в теории
универсальной аппроксимации, являются пространствами Фреше, то есть
полными метризуемыми локально выпуклыми топологическими векторными
пространствами. К ним относятся, в частности,
$C(\Omega,\mathbb{R}^n)$ с топологией равномерной сходимости на компактах,
$L^p_\mu(\mathbb{R}^m,\mathbb{R}^n)$ и липшицевы свободные пространства.
Важным свойством таких пространств является бэровость, что делает применимыми
результаты топологической динамики, в частности теорему Биркгофа о транзитивности,
используемую в дальнейшем для анализа универсального аппроксимационного свойства.

Пусть $X$ --- топологическое пространство функций (например,
$C(\Omega,\mathbb{R}^n)$ с топологией $d_{\mathrm{ucc}}$), а $\mathcal{G}\subset X$ ---
некоторый класс функций.

\begin{definition}[УАС относительно топологии]
Класс функций $\mathcal{G}$ обладает универсальным аппроксимационным свойством
в пространстве $X$, если
\[
\overline{\mathcal{G}}^X = X,
\]
то есть $\mathcal{G}$ плотно в $X$ относительно выбранной топологии.
\end{definition}

В классических теоремах универсальной аппроксимации рассматривалось пространство
$C(K,\mathbb{R})$ с нормой $\|\cdot\|_\infty$ на компакте $K$.
В мета-теоретическом подходе Крациоса подчёркивается, что выбор топологии является
существенным и напрямую влияет на корректность и форму утверждений об УАС.

\subsection{Архитектура в абстрактном смысле}
Для описания различных классов моделей вводится абстрактное понятие архитектуры.

\begin{definition}[Архитектура]
Архитектурой на функциональном пространстве $X$ называется пара
$\mathcal{A}=(\mathcal{F},\circ)$, где
\begin{itemize}
  \item $\mathcal{F}$ --- множество элементарных (примитивных) функций,
  \item $\circ$ --- правило композиции, позволяющее из конечного набора элементов
  $\mathcal{F}$ строить новые функции из $X$.
\end{itemize}
\end{definition}
\noindent
Множество всех функций, реализуемых с помощью архитектуры
$\mathcal{A}=(\mathcal{F},\circ)$, обозначается
\[
\mathcal{NN}((\mathcal{F},\circ)) \coloneq \mathcal{NN}(\mathcal{F},\circ)\subset X
\]
и называется \emph{множеством реализаций} архитектуры.
Архитектура обладает УАС в $X$, если $\mathcal{NN}(\mathcal{F},\circ)$ плотно в $X$.
Данное определение охватывает полносвязные нейронные сети, сверточные сети,
остаточные архитектуры, а также операторные сети вида
$\mathcal{N} = R\circ A \circ E$.

\begin{example}
  Абстрактное определение архитектуры охватывает широкий класс моделей,
  используемых в машинном обучении.
  \begin{itemize}
    \item \emph{Глубокие полносвязные сети:}
    $\mathcal{F}$ состоит из аффинных отображений, а оператор $\circ$
    реализует их композицию с поэлементным применением функции активации.
    \item \emph{Сверточные и структурированные сети:}
    $\mathcal{F}$ содержит отображения с жёстко заданной структурой
    (например, свёрточные или тоеплицевы операторы), а правило $\circ$
    допускает параллельные и последовательные композиции.
    \item \emph{Операторные сети:}
    архитектуры вида $N = R \circ A \circ E$, где $E$ --- кодировщик входных функций,
    $A$ --- нелинейный оператор в скрытом пространстве, а $R$ --- декодер.
  \end{itemize}
  
\end{example}

Ключевая идея мета-теории состоит в сведении задачи универсальной аппроксимации
к вопросам топологической динамики.
Пусть $X$ --- топологическое пространство и $T\colon X\to X$ --- непрерывное отображение.

\begin{definition}[Топологическая транзитивность]
Отображение $T$ называется топологически транзитивным, если для любых непустых
открытых множеств $U,V\subset X$ существует $n\in\mathbb{N}$ такое, что
$T^n(U)\cap V\neq\varnothing$.
\end{definition}

\begin{definition}[Плотная орбита]
Точка $x\in X$ имеет плотную орбиту, если множество
$\{T^n(x)\mid n\in\mathbb{N}\}$ плотно в $X$.
Такие точки называются гиперциклическими.
\end{definition}
\noindent
Интуитивно, наличие плотных орбит позволяет порождать плотные подмножества
пространства $X$, начиная с одного <<семени>> и многократно применяя оператор $T$.
Теорема Биркгофа о транзитивности утверждает, что в Бэровских пространствах
топологическая транзитивность эквивалентна существованию плотных орбит
\cite{Kratsios2021}.
Именно этот динамический механизм лежит в основе обобщённой теории УАС.

\begin{remark}
  В мета-теории Крациоса универсальное аппроксимационное свойство исследуется через четыре фундаментальных вопроса:
  \begin{enumerate}
    \item \emph{Характеризация:} при каких условиях архитектура обладает УАС?
    \item \emph{Конструкция:} как строить новые универсальные архитектуры из известных?
    \item \emph{Представление:} как связаны между собой все универсальные архитектуры?
    \item \emph{Существование:} существуют ли <<малые>> универсальные аппроксиматоры?
  \end{enumerate}
\end{remark}


\subsection{Характеризация универсального аппроксимационного свойства через транзитивную динамику}
Пусть $\mathcal{X}$ --- функциональное пространство, рассматриваемое как
топологическое пространство (как правило, сепарабельное, бесконечномерное
пространство Фреше или локально выпуклое), оснащённое метрикой $d_{\mathcal{X} }$.
Типичными примерами являются пространства $C(\Omega,\mathbb{R}^n)$ с топологией
равномерной сходимости на компактах.

Пусть $(\mathcal{F},\circ)$ --- архитектура в смысле определения предыдущего
раздела, а
\[
\mathcal{NN}(\mathcal{F},\circ) \subset \mathcal{X}
\]
--- множество всех функций, реализуемых данной архитектурой.
Универсальное аппроксимационное свойство (УАС) формулируется как условие плотности
$\mathcal{NN}(\mathcal{F},\circ)$ в $\mathcal{X}$.

\begin{theorem}[Крациос, динамическая характеризация УАС {\cite{Kratsios2021}}]
\label{thm:kratsios_characterization}
Пусть $\mathcal{X}$ --- топологическое пространство, гомеоморфное
бесконечномерному пространству Фреше.
Пусть $(\mathcal{F},\circ)$ --- архитектура на $\mathcal{X}$.
Тогда архитектура $(\mathcal{F},\circ)$ обладает универсальным аппроксимационным
свойством в $\mathcal{X}$, то есть $\overline{\mathcal{NN}(\mathcal{F},\circ)}^{\,\mathcal{X}} = \mathcal{X}$
тогда и только тогда, когда существуют семейство подпространств $\{X_i\}_{i\in I}$, $X_i \subset \mathcal{X}$,
непрерывные отображения $\varphi_i\colon X_i \to X_i$ и элементы $g_i \in \mathcal{NN}(\mathcal{F},\circ) \cap X_i$, такие, что:
\begin{enumerate}[label=(\alph*)]
  \item $\bigcup_{i\in I} X_i$ плотно в $\mathcal{X}$;
  \item для каждого $i\in I$ отображение $\varphi_i$ топологически
  транзитивно на $X_i$;
  \item для каждого $i\in I$ точка $g_i$ является гиперциклической для
  $\varphi_i$, причём
  \[
  \mathrm{Orb}_{\varphi_i}(g_i)
  \subset
  \mathcal{NN}(\mathcal{F},\circ) \cap X_i;
  \]
  \item[(d)] каждое пространство $X_i$ гомеоморфно пространству $C(\mathbb{R})$.
\end{enumerate}
Если пространство $\mathcal{X}$ является сепарабельным, то множество индексов
$I$ может быть выбрано одноточечным. $I = \{*\}$
\end{theorem}

\begin{remark}
Теорема утверждает, что универсальное аппроксимационное свойство эквивалентно
существованию \emph{реализуемых плотных орбит} подходящей динамики на
функциональном пространстве. Тем самым глубина, композиция и повторное применение
одних и тех же структурных блоков интерпретируются как итерации динамической
системы, а универсальность --- как плотность достижимых состояний.
\end{remark}

\begin{remark}
  В \cite{Kratsios2021} теорема выводится из более общего
  утверждения (Лемма~2 в Приложении~A), где эквивалентности формулируются в нескольких
  технических версиях. Ниже приводится доказательство в форме, соответствующей
  Лемме~2, из которого следует и теорема.
\end{remark}

\begin{lemma}[Характеризация универсального аппроксимационного свойства {\cite[Лемма~2]{Kratsios2021}}]
\label{lem:kratsios_lemma2_exact}
Пусть $\mathcal{X}$ --- функциональное пространство, $E$ --- бесконечномерное
пространство Фреше, и существует гомеоморфизм $\varphi:\mathcal{X}\to E$.
Пусть $(\mathcal{F},\circ)$ --- архитектура на $\mathcal{X}$.
Тогда следующие утверждения эквивалентны:

\begin{enumerate}[label=(\roman*)]
\item \textbf{(УАС).}
Архитектура $(\mathcal{F},\circ)$ обладает универсальным аппроксимационным свойством на $\mathcal{X}$, то есть
\[
\overline{\mathcal{NN}(\mathcal{F},\circ)}^{\,\mathcal{X}}=\mathcal{X}.
\]

\item \textbf{(Декомпозиция УАС через подпространства).}
Существует семейство подпространств $\{X_i\}_{i\in I}$ пространства $\mathcal{X}$ такое, что:
\begin{enumerate}[label=(\alph*)]
  \item $\displaystyle \bigcup_{i\in I}X_i$ плотно в $\mathcal{X}$;
  \item для каждого $i\in I$ множество $\varphi(X_i)$ является
  сепарабельным бесконечномерным фреше-подпространством $E$, и
  $\mathcal{NN}(\mathcal{F},\circ)\cap X_i$ содержит
  счётное плотное линейно независимое подмножество $\varphi(X_i)$;
  \item для каждого $i\in I$ существует гомеоморфизм
  \[
  \phi_i\colon X_i \to L^2(\mathbb{R}).
  \]
\end{enumerate}

\item\textbf{(Декомпозиция УАС через топологически транзитивную динамику).}
Существует семейство подпространств $\{X_i\}_{i\in I}$ пространства $\mathcal{X}$
и непрерывных отображений $\{\phi_i\}_{i\in I}$, где $\phi_i\colon X_i\to X_i$, такое, что:
\begin{enumerate}[label=(\alph*)]
  \item $\displaystyle \bigcup_{i\in I}X_i$ плотно в $\mathcal{X}$;
  \item для каждого $i\in I$ и любых непустых открытых $U,V\subset X_i$
  существует $N_{i,U,V}\in\mathbb{N}$ такое, что
  \[
  \phi_i^{N_{i,U,V}}(U)\cap V\neq\varnothing;
  \]
  \item для каждого $i\in I$ существует $g_i\in \mathcal{NN}(\mathcal{F},\circ)\cap X_i$
  такое, что множество $\{\phi_i^n(g_i)\}_{n\in\mathbb{N}}$ плотно в
  $\mathcal{NN}(\mathcal{F},\circ)\cap X_i$ и, в частности, плотно в $X_i$;
  \item для каждого $i\in I$ пространство $X_i$ гомеоморфно $C(\mathbb{R})$.
\end{enumerate}

\item \textbf{(Параметризация УАС на подпространствах).}
Существует семейство троек $\{(X_i,\phi_i,\psi_i)\}_{i\in I}$, где
$X_i$ --- сепарабельные топологические пространства,
$\phi_i\colon X_i\to X_i$ --- непрерывные неконстантные отображения и
$\psi_i\colon X_i\to\mathcal{X}$ --- непрерывные отображения, такое что:
\begin{enumerate}[label=(\alph*)]
  \item $\displaystyle \bigcup_{i\in I}\psi_i(X_i)$ плотно в $\mathcal{X}$;
  \item для каждого $i\in I$ и любых непустых открытых $U,V\subset \mathcal{X}_i$
  существует $N_{i,U,V}\in\mathbb{N}$ такое, что
  \[
  \phi_i^{N_{i,U,V}}(U)\cap V\neq\varnothing;
  \]
  \item для каждого $i\in I$ существует $x_i\in \mathcal{NN}(\mathcal{F},\circ)\cap \psi_i(X_i)$
  такое, что множество
  \[
  \{\psi_i\circ \phi_i^n(x_i)\}_{n\in\mathbb{N}}
  \]
  плотно в $\mathcal{NN}(\mathcal{F},\circ)\cap \psi_i(X_i)$ и, в частности, плотно в $\psi_i(X_i)$.
\end{enumerate}
\end{enumerate}
Более того, если $\mathcal{X}$ сепарабельно, то индексное множество $I$ можно выбрать одноточечным.
\end{lemma}

\begin{proof}
Обозначим $\mathcal{N} \coloneq \mathcal{NN}(\mathcal{F},\circ) \subset X$.

\noindent\textbf{Шаг 1. (ii) $\Rightarrow$ (i): <<плотность на подпространствах $\Rightarrow$ УАС>>.}
Предположим, что существует семейство подпространств $\{X_i\}_{i\in I}$ такое, что
$\bigcup_{i\in I} X_i$ плотно в $X$, и при этом $\mathcal{N}\cap X_i$ плотно в $X_i$ для каждого $i$.
Рассмотрим объединение
\[
\widetilde{X} \coloneq \bigcup_{i\in I} X_i,
\]
оснащённое \emph{наитончайшей топологией}, делающей каждое $X_i$ подпространством
(в статье это фиксируется ссылкой на стандартный факт о финальной топологии на объединении
подпространств).

Тогда достаточно показать плотность $\bigcup_{i\in I}(\mathcal{N}\cap X_i)$ в $\widetilde{X}$.
Пусть $U\subset \widetilde{X}$ --- непустое открытое множество. По определению топологии на
$\widetilde{X}$ существует $i\in I$, для которого $U\cap X_i$ непусто и открыто в $X_i$.
Так как $\mathcal{N}\cap X_i$ плотно в $X_i$, имеем
\[
(U\cap X_i)\cap(\mathcal{N}\cap X_i)\neq\varnothing,
\]
откуда $U\cap \mathcal{N}\neq\varnothing$. Значит, $\mathcal{N}$ плотно в $\widetilde{X}$, а так как
$\widetilde{X}$ плотно в $\mathcal{X}$, получаем плотность $\mathcal{N}$ в $\mathcal{X}$, то есть УАС.

\noindent\textbf{Шаг 2. (iii) $\Rightarrow$ (i): <<реализуемая плотная орбита $\Rightarrow$ УАС>>.}
Предположим, что для каждого $i\in I$ существует непрерывное отображение
$\varphi_i:X_i\to X_i$ и элемент $g_i\in \mathcal{N}\cap X_i$ такой, что орбита
$\{\varphi_i^n(g_i)\}_{n\in\mathbb{N}}$ плотна в $X_i$ и целиком лежит в $\mathcal{N}\cap X_i$.
Тогда $\mathcal{N}\cap X_i$ содержит плотное подмножество $X_i$, следовательно,
$\mathcal{N}\cap X_i$ плотно в $X_i$ для всех $i$. Поэтому условия пункта (ii) выполнены,
и по Шагу 1 имеем УАС в $\mathcal{X}$.

\noindent\textbf{Шаг 3. (i) $\Rightarrow$ (ii): <<УАС $\Rightarrow$ разложение на сепарабельные фреше-подпространства>>.}
Пусть $\mathcal{N}$ плотно в $\mathcal{X}$ и $\mathcal{X}$ гомеоморфно бесконечномерному фреше-пространству $E$
(то есть существует гомеоморфизм $\Phi\colon X\to E$).
Используя стандартную структуру фреше-пространств, в статье выбирают семейство
\[
\{X_i\}_{i\in I}\quad\text{сепарабельных бесконечномерных фреше-подпространств }X_i\subset \mathcal{X},
\]
такое что $\bigcup_{i\in I}X_i$ плотно в $\mathcal{X}$ и, кроме того, пересечение $\mathcal{N}\cap X_i$
плотно в каждом $X_i$.

Далее, так как $X_i$ сепарабельно и бесконечномерно, можно выбрать в $\mathcal{N}\cap X_i$
\emph{счётное плотное линейно независимое} множество
\[
\{d_{n,i}\}_{n\in\mathbb{N}}\subset \mathcal{N}\cap X_i,
\]
что и составляет усиление, требуемое в (ii.b).
Наконец, для каждого $i$ фиксируется гомеоморфизм $X_i\simeq L^2(\mathbb{R})$ (ii.c).

\noindent\textbf{Шаг 4. (ii) $\Rightarrow$ (iii): построение динамики со сдвигом плотного независимого набора.}
Предположим, что выполнено (ii). Зафиксируем $i\in I$ и положим
\[
E_i \coloneq \mathrm{span}\{d_{n,i}:n\in\mathbb{N}\}.
\]
Тогда $E_i$ --- счётномерное линейное подпространство, плотное в $X_i$.

Определим на $E_i$ линейный оператор-сдвиг $T_i\colon E_i\to E_i$ по правилу
\[
T_i(d_{n,i}) \coloneq d_{n+1,i},\qquad n\in\mathbb{N},
\]
и продолжим его линейно на весь $E_i$. По построению получаем
\[
T_i^n(d_{0,i}) = d_{n,i},\qquad n\in\mathbb{N}.
\]
Следовательно, орбита $\{T_i^n(d_{0,i})\}_{n\in\mathbb{N}}$ плотна в $E_i$, а значит
плотна и в $X_i$.

Далее в статье применяется стандартный результат из линейной динамики на фреше-пространствах:
оператор $T_i$, заданный на плотном счётномерном подпространстве, можно продолжить
до \emph{непрерывного линейного} оператора $\widehat{T}_i\colon X_i\to X_i$, сохраняющего соотношение
$\widehat{T}_i^n(d_{0,i})=d_{n,i}$ для всех $n$. Обозначим
\[
\varphi_i \coloneq \widehat{T}_i,\qquad g_i \coloneq d_{0,i}.
\]
Тогда $\{\varphi_i^n(g_i)\}_{n\in\mathbb{N}}=\{d_{n,i}\}_{n\in\mathbb{N}}$ плотна в $X_i$.
Кроме того, по выбору $\{d_{n,i}\}\subset \mathcal{N}\cap X_i$ имеем
\[
\mathrm{Orb}_{\varphi_i}(g_i) \subset \mathcal{N}\cap X_i,
\]
то есть выполняется требование <<реализуемости орбиты>> (iii.c).

Осталось получить (iii.b), то есть топологическую транзитивность $\varphi_i$ на $X_i$.
В статье это делается через теорему Биркгофа о транзитивности: так как $X_i$
гомеоморфно $L^2(\mathbb{R})$ (полное сепарабельное метрическое пространство без
изолированных точек), а $\varphi_i$ имеет гиперциклическую точку (поскольку
орбита $g_i$ плотна), то по теореме Биркгофа $\varphi_i$ топологически транзитивно.
Тем самым (iii.b) доказано.

\noindent\textbf{Шаг 5. Почему $X_i\simeq C(\mathbb{R})$.}
В статье пункт (iii.d) получается так: $X_i\simeq L^2(\mathbb{R})$ по (ii.c), а затем
используется теорема Андерсона--Кадецa о том, что любые два сепарабельных
бесконечномерных фреше-пространства гомеоморфны; в частности,
$L^2(\mathbb{R}) \simeq C(\mathbb{R})$. Следовательно, $X_i\simeq C(\mathbb{R})$.

\noindent\textbf{Шаг 6. (iii) $\Rightarrow$ (iv)}
Достаточно для каждого $i \in I$ взять $\phi_i \coloneq id_{X_i}$, $\phi_i \coloneq \phi_i$, $x_i = g_i$. Это приводит к (iv).

\noindent\textbf{Шаг 7. (iv) $\Rightarrow$ (i)}
По (iv.c) $\mathcal{N} \cap X_i$ плотно в $X_i$, поэтому 
\[
\bigcup_{i \in I} X_i = \bigcup_{i \in I} \overline{\mathcal{N} \cap X_i} \subset \overline{\bigcup_{i \in I} \mathcal{N} \cap X_i} \subset \mathcal{X}.
\]
По (iv.a) $\bigcup_{i \in I} X_i$ плотно в $\mathcal{X}$, а значит минимальным замкнутым множеством, содержащим это объединение, будет всё пространство.
По включению выше, наименьшим множеством, содержащим $bigcup_{i \in I} \mathcal{N} \cap X_i$ должно быть $\mathcal{X}$.
Поэтому $\mathcal{N}$ плотно в $\mathcal{X}$, и справедливо (i).

Таким образом, доказаны импликации (ii)$\Rightarrow$(i), (iii)$\Rightarrow$(i),
(i)$\Rightarrow$(ii), (ii)$\Rightarrow$(iii), (ii)$\Rightarrow$(iv) и (iv)$\Rightarrow$(i), что завершает доказательство эквивалентностей
и тем самым доказательство теоремы~\ref{thm:kratsios_characterization}.
\end{proof}
\begin{remark}
  Если само пространство $\mathcal{X}$ сепарабельно, то можно положить $I=\{\ast\}$ и взять $\mathcal{X}_\ast=\mathcal{X}$,
  что прямо указано в формулировке леммы.
\end{remark}

\subsection{Перенос универсального аппроксимационного свойства между пространствами}
Пусть $(\mathcal{X},d_\mathcal{X})$ и $(\mathcal{Y},d_\mathcal{Y})$ --- функциональные пространства, рассматриваемые как
топологические пространства (как правило, фреше-пространства или локально
выпуклые топологические векторные пространства).
Для подмножества $A$ в пространстве $Z$ обозначим через $\overline{A}^{\,Z}$
его замыкание в топологии пространства $Z$.

Пусть $(\mathcal{F},\circ)$ --- архитектура на пространстве $X$, а
\[
\mathcal{NN}(\mathcal{F},\circ)\subset \mathcal{X}
\]
--- соответствующее множество реализуемых функций.
Архитектура обладает универсальным аппроксимационным свойством (УАС) в $\mathcal{X}$, если
\[
\overline{\mathcal{NN}(\mathcal{F},\circ)}^{\,\mathcal{X}}=\mathcal{X}.
\]


\begin{theorem}[Крациос, принцип переноса УАС {\cite{Kratsios2021}}]
\label{thm:transport_uap}
Пусть $(\mathcal{F},\circ)$ --- архитектура на функциональном пространстве $\mathcal{X}$.
Предположим, что выполнены следующие условия:
\begin{enumerate}
\item[(A1)]
Архитектура $(\mathcal{F},\circ)$ обладает универсальным аппроксимационным
свойством в $\mathcal{X}$, то есть $\overline{\mathcal{NN}(\mathcal{F},\circ)}^{\,\mathcal{X}}=\mathcal{X}$.

\item[(A2)]
Существует семейство непрерывных отображений $\{\Psi_i\colon \mathcal{X}\to \mathcal{Y}\}_{i\in I}$,
такое, что объединение их образов плотно в $\mathcal{Y}$: $\overline{\bigcup_{i\in I}\Psi_i(\mathcal{X})}^{\,\mathcal{Y}}=\mathcal{Y}$.

\item[(A3)]
Для каждого $i\in I$ существует архитектура $(\mathcal{F},\circ)_i$ на $\mathcal{Y}$,
удовлетворяющая включению
\[
\Psi_i\bigl(\mathcal{NN}(\mathcal{F},\circ)\bigr)
\subset
\mathcal{NN}\bigl((\mathcal{F},\circ)_i\bigr)
\subset Y.
\]
\end{enumerate}

Тогда объединённая архитектура
\[
(\mathcal{F},\circ)_{\mathrm{tr}}
\coloneq
\bigsqcup_{i\in I}(\mathcal{F},\circ)_i,
\]
допускающая выбор индекса $i$ и последующее применение архитектуры
$(\mathcal{F},\circ)_i$, обладает универсальным аппроксимационным свойством в $Y$: $\overline{\mathcal{NN}\bigl((\mathcal{F},\circ)_{\mathrm{tr}}\bigr)}^{\,\mathcal{Y}} = \mathcal{Y}$.

Если пространство $\mathcal{Y}$ является сепарабельным, то индексное множество $I$ можно
выбрать счётным.
\end{theorem}


\begin{proof}[]
Обозначим $\mathcal{N}\coloneq\mathcal{NN}(\mathcal{F},\circ)\subset\mathcal{X}$ и
$\mathcal{N}_{\mathrm{tr}}\coloneq\mathcal{NN}\bigl((\mathcal{F},\circ)_{\mathrm{tr}}\bigr)\subset\mathcal{Y}$.
Достаточно показать, что каждое непустое открытое множество $V\subset\mathcal{Y}$
пересекает $\mathcal{N}_{\mathrm{tr}}$.

Пусть $V\subset\mathcal{Y}$ --- непустое открытое. Из (A2) следует, что
$\bigcup_{i\in I}\Psi_i(\mathcal{X})$ плотно в $\mathcal{Y}$, значит
\[
V\cap \bigcup_{i\in I}\Psi_i(\mathcal{X})\neq\varnothing.
\]
Следовательно, существуют $i\in I$ и $x\in\mathcal{X}$ такие, что $\Psi_i(x)\in V$.
По непрерывности $\Psi_i$ существует открытая окрестность $U\subset\mathcal{X}$ точки $x$,
для которой
\[
\Psi_i(U)\subset V.
\]

Далее, из (A1) имеем плотность $\mathcal{N}$ в $\mathcal{X}$, значит $U\cap\mathcal{N}\neq\varnothing$.
Выберем $h\in U\cap\mathcal{N}$. Тогда $\Psi_i(h)\in \Psi_i(U)\subset V$.

Наконец, по (A3) имеем $\Psi_i(h)\in \Psi_i(\mathcal{N})\subset \mathcal{NN}((\mathcal{F},\circ)_i)$.
Так как $(\mathcal{F},\circ)_{\mathrm{tr}}$ допускает выбор ветви $i$, то
\[
\mathcal{NN}((\mathcal{F},\circ)_i)\subset \mathcal{N}_{\mathrm{tr}}.
\]
Следовательно, $\Psi_i(h)\in V\cap \mathcal{N}_{\mathrm{tr}}$, то есть каждое непустое открытое $V$
пересекает $\mathcal{N}_{\mathrm{tr}}$. Это эквивалентно плотности $\mathcal{N}_{\mathrm{tr}}$ в $\mathcal{Y}$.
\end{proof}

Отдельно рассмотрим два важных частных случая.\\
\textbf{Случай 1: один гомеоморфизм.}
Пусть $\Psi\colon \mathcal{X}\to\mathcal{Y}$ --- гомеоморфизм. Если $\mathcal{N}\subset\mathcal{X}$ плотно в $\mathcal{X}$,
то $\Psi(\mathcal{N})$ плотно в $\mathcal{Y}$, то есть
\[
\overline{\Psi(\mathcal{N})}^{\,\mathcal{Y}}=\mathcal{Y}.
\]
В частности, если архитектура $(\mathcal{F},\circ)$ имеет УАС в $\mathcal{X}$, то любой класс в $\mathcal{Y}$,
содержащий $\Psi(\mathcal{NN}(\mathcal{F},\circ))$, является универсальным в $\mathcal{Y}$.
\begin{proof}[]
Пусть $V\subset\mathcal{Y}$ непусто и открыто. Тогда $\Psi^{-1}(V)$ непусто и открыто в $\mathcal{X}$
(так как $\Psi$ --- гомеоморфизм). Плотность $\mathcal{N}$ в $\mathcal{X}$ даёт
$\Psi^{-1}(V)\cap \mathcal{N}\neq\varnothing$. Применяя $\Psi$, получаем
$V\cap \Psi(\mathcal{N})\neq\varnothing$. Следовательно, $\Psi(\mathcal{N})$ плотно в $\mathcal{Y}$.
\end{proof}

\textbf{Случай 2: плотное семейство вложений}
Пусть $\{\Psi_i\colon \mathcal{X}\to\mathcal{Y}\}_{i\in I}$ --- семейство непрерывных отображений, причём каждое $\Psi_i$
является вложением (homeomorphism onto its image), и
\[
\overline{\bigcup_{i\in I}\Psi_i(\mathcal{X})}^{\,\mathcal{Y}}=\mathcal{Y}.
\]
Если $(\mathcal{F},\circ)$ имеет УАС в $\mathcal{X}$ и для каждого $i$ существует архитектура
$(\mathcal{F},\circ)_i$ на $\mathcal{Y}$, содержащая $\Psi_i(\mathcal{NN}(\mathcal{F},\circ))$,
то объединённая архитектура $(\mathcal{F},\circ)_{\mathrm{tr}}=\bigsqcup_{i\in I}(\mathcal{F},\circ)_i$
имеет УАС в $\mathcal{Y}$.
\begin{proof}
Это прямое применение принципа переноса: условие плотного покрытия образами $\Psi_i(\mathcal{X})$
есть (A2), а включения реализуемостей есть (A3). Свойство <<вложение>> здесь даёт удобную
интерпретацию: на каждом $\Psi_i(\mathcal{X})$ можно аппроксимировать после
обратного перехода в $\mathcal{X}$, а плотность покрытия поднимает результат на всё $\mathcal{Y}$.
Формально доказательство совпадает с доказательством теоремы выше.
\end{proof}

В предыдущем разделе был сформулирован общий принцип переноса универсального
аппроксимационного свойства между функциональными пространствами, основанный на
непрерывных отображениях и плотности их образов.
Данный принцип носит архитектурно-агностичный характер и позволяет переносить
универсальность с одного пространства функций на другое при минимальных
топологических предположениях.

В статье \cite{Kratsios2021} этот общий подход реализуется в более конкретной и
конструктивной форме, ориентированной на практическое построение новых
универсальных архитектур. В качестве исходного пространства выступает
пространство непрерывных функций
\[
\mathcal{X}_0 \coloneq C(\mathbb{R}^m,\mathbb{R}^n),
\]
на котором уже задана архитектура, обладающая универсальным аппроксимационным
свойством. Новые универсальные аппроксиматоры в целевом функциональном
пространстве $\mathcal{X}$ строятся посредством применения семейства непрерывных
преобразований, отображающих $\mathcal{X}_0$ в $\mathcal{X}$.

Следующая теорема, являющаяся теоремой~2 в \cite{Kratsios2021}, формализует этот
конструктивный принцип и показывает, что любой универсальный аппроксиматор на
$C(\mathbb{R}^m,\mathbb{R}^n)$ порождает универсальный аппроксиматор в более общем
функциональном пространстве $\mathcal{X}$ посредством семейства непрерывных
преобразований.

\begin{theorem}[Крациос, универсальные аппроксиматоры через преобразования {\cite{Kratsios2021}}]
\label{thm:universal_by_transformation}
Пусть $m,n\in\mathbb{N}_+$,
$\mathcal{X}$ --- функциональное пространство,
и пусть $(\mathcal{F},\circ)$ является универсальным аппроксиматором
в пространстве
\[
\mathcal{X}_0 \coloneq C(\mathbb{R}^m,\mathbb{R}^n).
\]

Пусть $\{\Phi_i\}_{i\in I}$ --- непустое семейство непрерывных отображений
\[
\Phi_i : \mathcal{X}_0 \longrightarrow \mathcal{X},
\]
удовлетворяющее условию плотности
\begin{equation}
\label{eq:dense_union_transform}
\overline{
\bigcup_{i\in I}
\Phi_i\bigl(\mathcal{X}_0\bigr)
}^{\,\mathcal{X}}
=
\mathcal{X}.
\end{equation}

Тогда архитектура $(\mathcal{F}_{\Phi},\circ_{\Phi})$ обладает универсальным
аппроксимационным свойством в пространстве $\mathcal{X}$, где:
\begin{itemize}
\item $\mathcal{F}_{\Phi} \coloneq \mathcal{F} \times I$;
\item правило композиции $\circ_{\Phi}$ задаётся формулой
\[
\circ_{\Phi}\bigl((f_j,i_j)_{j=1}^J\bigr)
\coloneq
\Phi_{i_J}\!\left(
\circ\bigl((f_j)_{j=1}^J\bigr)
\right).
\]
\end{itemize}
\end{theorem}

\begin{proof}
По теореме Андерсона--Кадца \cite[Theorem~73]{AndersonKadec}
пространства $C(\mathbb{R}^m,\mathbb{R}^n)$ и $C(\mathbb{R})$ гомеоморфны.
Следовательно, без ограничения общности можно считать, что $m=n=1$.

Обозначим
\[
\mathcal{X}' \coloneq \bigcup_{i\in I} \Phi_i\bigl(C(\mathbb{R})\bigr).
\]
По предположению теоремы $\mathcal{X}'$ плотно в $\mathcal{X}$.
Так как плотность является транзитивным свойством, достаточно показать, что
\[
\bigcup_{i\in I}\mathcal{NN}(\mathcal{F},\circ)\cap \Phi_i\bigl(C(\mathbb{R})\bigr)
\]
плотно в $\mathcal{X}'$; тогда из плотности $\mathcal{X}'$ в $\mathcal{X}$ будет
следовать плотность этого множества и в $\mathcal{X}$.

Поскольку каждое отображение $\Phi_i$ непрерывно, топология,
индуцированная на $\Phi_i(C(\mathbb{R}))$ как подпространстве $\mathcal{X}$,
является не тоньше финальной топологии, относительно которой каждое
$\Phi_i(C(\mathbb{R}))$ является подпространством.
Обозначим через
\[
\mathcal{X}'' \coloneq \bigsqcup_{i\in I}\Phi_i\bigl(C(\mathbb{R})\bigr)
\]
дизъюнктное объединение образов, снабжённое наитончайшей топологией,
делающей каждое вложение $\Phi_i(C(\mathbb{R})) \hookrightarrow \mathcal{X}''$
непрерывным.

По построению, если $U\subset \mathcal{X}'$ открыто, то $U$ открыто и в
$\mathcal{X}''$, и наоборот.
Следовательно, для доказательства плотности достаточно показать, что
\[
\bigcup_{i\in I}\mathcal{NN}(\mathcal{F},\circ)\cap \Phi_i\bigl(C(\mathbb{R})\bigr)
\]
плотно в $\mathcal{X}''$.

Далее, пространство $\mathcal{X}''$ может быть представлено как
топологическое фактор-пространство дизъюнктного объединения
$\bigsqcup_{i\in I} C(\mathbb{R})$
по эквивалентности
\[
f_i \sim f_j
\quad\Longleftrightarrow\quad
\Phi_i(f_i)=\Phi_j(f_j).
\]
Обозначим соответствующее фактор-отображение через
\[
Q_{\mathcal{X}''}\colon \bigsqcup_{i\in I} C(\mathbb{R}) \longrightarrow \mathcal{X}''.
\]

Пусть $U\subset\mathcal{X}''$ --- непустое открытое множество.
Тогда прообраз $Q_{\mathcal{X}''}^{-1}(U)$ открыт в дизъюнктном объединении,
а значит существует $i\in I$ и открытое множество
$U'\subset C(\mathbb{R})$ такое, что
\[
\varnothing \neq Q_{\mathcal{X}''}(U') \subset U \cap \Phi_i\bigl(C(\mathbb{R})\bigr).
\]

Поскольку архитектура $(\mathcal{F},\circ)$ является универсальным
аппроксиматором на $C(\mathbb{R})$, множество
$\mathcal{NN}(\mathcal{F},\circ)\cap C(\mathbb{R})$
плотно в $C(\mathbb{R})$.
Следовательно,
\[
U' \cap \mathcal{NN}(\mathcal{F},\circ)\neq\varnothing,
\]
и потому
\[
\varnothing \neq
Q_{\mathcal{X}''}\!\left(
U' \cap \mathcal{NN}(\mathcal{F},\circ)
\right)
\subset
U \cap
\mathcal{NN}(\mathcal{F},\circ)\cap \Phi_i\bigl(C(\mathbb{R})\bigr).
\]

Из этого следует, что для любого непустого открытого множества
$U\subset\mathcal{X}''$ выполняется
\[
\varnothing \neq
U \cap
\bigcup_{i\in I}
\mathcal{NN}(\mathcal{F},\circ)\cap \Phi_i\bigl(C(\mathbb{R})\bigr).
\]

Тем самым
$\bigcup_{i\in I}\mathcal{NN}(\mathcal{F},\circ)\cap \Phi_i(C(\mathbb{R}))$
плотно в $\mathcal{X}''$, а значит плотно и в $\mathcal{X}'$.
Так как $\mathcal{X}'$ плотно в $\mathcal{X}$, получаем плотность этого
множества и в $\mathcal{X}$.

Следовательно, архитектура $(\mathcal{F}_{\Phi},\circ_{\Phi})$
обладает универсальным аппроксимационным свойством в $\mathcal{X}$.
\end{proof}
\end{document}